\chapter{Braided Factorization}
\label{ch:braided-factorization}

%%: The E_1/ordered story is primitive. The E_2 braided
%% structure is the Drinfeld double / center construction applied to E_1.
%%: kappa always subscripted in Vol III.
%% Prerequisite: Chapter~\ref{ch:quantum-groups} (Quantum Groups: Foundations)
%% owns the Drinfeld--Jimbo presentation of $\Uq(\frakg)$, the universal
%% $R$-matrix, the quantum Yang--Baxter equation, and the Kazhdan--Lusztig
%% equivalence. The present chapter builds the $E_2$-chiral / factorization
%% layer on top.

A braided monoidal category is a monoidal category equipped with a
coherent system of isomorphisms $V \otimes W \xrightarrow{\sim} W \otimes V$
satisfying the hexagon axioms. The bar-cobar adjunction of Volume~I,
extended to the $E_2$ setting, produces factorization coalgebras whose
arity-$(1,1)$ component is a categorical $R$-matrix. Chapter~\ref{ch:quantum-groups}
collects the classical Drinfeld--Jimbo presentation of $\Uq(\frakg)$, the
universal $R$-matrix, and the quantum Yang--Baxter equation; here we assume
that material and construct its $E_2$-chiral / factorization avatar. The
fundamental problem solved in this chapter: how does the $E_2$-chiral bar
complex $B_{E_2}(\cA)$ encode a categorical braiding, and what is the precise
relationship between the braided Koszul dual and the quantum group targets
of Conjecture~\ref{conj:qgf-cy-c}?

\section{The categorical $R$-matrix and Yang--Baxter equation}
\label{sec:bf-categorical-r-matrix}

The Drinfeld--Jimbo universal $R$-matrix of
Proposition~\ref{prop:qgf-classical-limit-r} lives in $\Uq(\frakg)^{\hat\otimes 2}$
and is an \emph{algebraic} datum. The Vol~III programme produces $R$-matrices
of a different kind: \emph{categorical} $R$-matrices, extracted from the
arity-$(1,1)$ component of an $E_2$-chiral bar complex. The two pictures
are compatible through Conjecture~CY-C (Conjecture~\ref{conj:qgf-cy-c}):
when the CY-to-chiral functor produces an affine Kac--Moody chiral algebra,
the categorical $R$-matrix coincides with the Drinfeld--Jimbo $R$-matrix
under the Kazhdan--Lusztig equivalence (Theorem~\ref{thm:qgf-kazhdan-lusztig}).

\begin{remark}[Drinfeld--Jimbo vs.\ FRT, from the bar-complex standpoint]
\label{rem:bf-dj-vs-frt}
The Drinfeld--Jimbo presentation constructs $\Uq(\frakg)$ from generators
and relations; the Faddeev--Reshetikhin--Takhtajan (FRT) construction
recovers $\Uq(\frakg)$ from the $R$-matrix via the RTT relation
$R_{12} T_1 T_2 = T_2 T_1 R_{12}$. The FRT construction is the
bar-complex-natural one: the $R$-matrix lives in arity $(1,1)$ of
$B^{\mathrm{ord}}(A)$, and the RTT relation is the
arity-$(1,1,1)$ bar differential. This is the standpoint
of the present chapter: the categorical $R$-matrix below is the
arity-$(1,1)$ component of $B_{E_2}(\Phi(\cC))$, and the parametric
Yang--Baxter equation is coassociativity of the deconcatenation coproduct.
\end{remark}

\begin{definition}[Categorical $R$-matrix]
\label{def:categorical-r-matrix}
For a CY category $\cC$ of dimension $d = 2$, the \emph{categorical
$R$-matrix} $\cR_\cC(z)$ is the arity-$(1,1)$ component of the
$E_2$-bar complex $B_{E_2}(\Phi(\cC))$:
\[
 \cR_\cC(z) \in \End(V \otimes V)[[z, z^{-1}]]
\]
where $V = s^{-1}\bar{A}_\cC$ is the desuspended generating space and
$z$ is the spectral parameter. The assignment $\cC \mapsto \cR_\cC(z)$
is functorial in CY categories (for exact CY functors preserving the
$E_2$-structure).
\end{definition}

\begin{proposition}[Yang--Baxter from bar associativity]
\label{prop:ybe-from-bar}
\ClaimStatusProvedElsewhere
The categorical $R$-matrix $\cR_\cC(z)$ satisfies the Yang--Baxter
equation
\[
 \cR_{12}(z_1 - z_2) \, \cR_{13}(z_1 - z_3) \, \cR_{23}(z_2 - z_3)
 = \cR_{23}(z_2 - z_3) \, \cR_{13}(z_1 - z_3) \, \cR_{12}(z_1 - z_2).
\]
This is a formal consequence of the coassociativity of the bar coproduct
$\Delta_{\mathrm{decat}}$ applied to arity $3$ of $B^{\mathrm{ord}}(A)$.
\end{proposition}

\begin{proof}
The bar complex $B^{\mathrm{ord}}(A)$ is a coassociative coalgebra with
deconcatenation coproduct. The two compositions
$(\Delta \otimes \id) \circ \Delta$ and $(\id \otimes \Delta) \circ \Delta$
coincide on arity-$3$ elements; the $R$-matrix arises from the bar
differential restricted to arity $(1,1)$. Coassociativity at arity $3$
gives $\cR_{12}\cR_{13}\cR_{23} = \cR_{23}\cR_{13}\cR_{12}$, which is
the Yang--Baxter equation. The spectral parameters arise from the
$z$-dependence of the collision residue (so the $R$-matrix has one
fewer pole order than the OPE). See Volume~II, Chapter~11 for the
full derivation.
\end{proof}


\section{Braided monoidal categories and $E_2$-algebras}
\label{sec:braided-e2}

The connection between $E_2$-algebras and braided monoidal categories is
given by the Joyal--Street coherence theorem and its
$\infty$-categorical refinement.

\begin{proposition}[Dunn additivity for $E_2$]
\label{prop:dunn-e2}
\ClaimStatusProvedElsewhere
The $E_2$ operad satisfies $\Etwo \simeq \Eone \otimes_{E_0} \Eone$.
An $\Etwo$-algebra is an $\Eone$-algebra in $\Eone$-algebras. In
particular, an $\Etwo$-algebra has two compatible associative products
whose interchange law is the braiding.
\end{proposition}

\begin{corollary}
\label{cor:e2-braided-monoidal}
The representation category $\Rep^{E_2}(\cA)$ of an $E_2$-algebra
$\cA$ is braided monoidal. The braiding on $V \otimes W$ is given by
the action of the $R$-matrix: $\sigma_{V,W} = \cR_{V,W} \circ \tau$
where $\tau \colon V \otimes W \to W \otimes V$ is the transposition
and $\cR_{V,W}$ is the $R$-matrix evaluated on the pair $(V,W)$.
\end{corollary}


\section{Factorization homology with $E_2$-coefficients}
\label{sec:fh-e2}

Factorization homology $\int_M \cA$ of an $\En$-algebra $\cA$ over an
$n$-manifold $M$ (Ayala--Francis) generalizes Hochschild homology
($n = 1$, $M = S^1$) and chiral homology ($n = 2$, $M$ a surface).
The relevance to Vol~III is that CY$_2$ categories produce $\Etwo$-chiral
algebras (by CY-A at $d = 2$), and factorization homology of these
$\Etwo$-algebras over surfaces computes genus-$g$ conformal block spaces;
the shadow obstruction tower of Vol~I controls the moduli dependence.

\begin{definition}[Factorization homology]
\label{def:factorization-homology}
Let $\cA$ be an $\En$-algebra and $M$ an $n$-manifold. The
\emph{factorization homology} $\int_M \cA$ is the colimit
\[
 \int_M \cA \;=\; \colim_{\mathrm{Disk}_{n/M}} \cA
\]
over the $\infty$-category $\mathrm{Disk}_{n/M}$ of embeddings of
disjoint unions of disks into $M$. The colimit is computed in the
$\infty$-category of chain complexes (Ayala--Francis, 2015).
\end{definition}

\begin{proposition}[Excision for factorization homology]
\label{prop:fh-excision}
\ClaimStatusProvedElsewhere
Let $M = M_1 \cup_{N \times \R} M_2$ be a collar-gluing decomposition
of an $n$-manifold. For an $\En$-algebra $\cA$, there is a canonical
equivalence
\[
 \int_M \cA
 \;\simeq\;
 \int_{M_1} \cA \;\otimes_{\int_{N \times \R} \cA}\; \int_{M_2} \cA,
\]
where the tensor product is taken over the $\Eone$-algebra
$\int_{N \times \R} \cA$ (Ayala--Francis excision).
\end{proposition}

Excision is the engine that converts local data (disk-level OPE) into
global invariants (genus-$g$ amplitudes). Applied to the pair-of-pants
decomposition $\Sigma_g = \#_{2g-2}(\text{pair of pants})$, excision
expresses $\int_{\Sigma_g} \cA$ as an iterated relative tensor product
over Hochschild homology, recovering the sewing prescription of conformal
field theory.

\begin{proposition}[$E_2$ factorization homology on surfaces]
\label{prop:fh-e2-surface}
\ClaimStatusProvedElsewhere
For an $\Etwo$-algebra $\cA$ and a closed oriented surface
$\Sigma_g$ of genus $g$:
\begin{enumerate}[label=(\roman*)]
 \item $\int_{\Sigma_g} \cA$ is the genus-$g$ conformal block space;
 \item $\int_{S^1 \times [0,1]} \cA \simeq \HH_\bullet(\cA)$, the
 Hochschild homology;
 \item $\int_{T^2} \cA \simeq \HH_\bullet(\HH_\bullet(\cA))$
 (double Hochschild homology, computing the genus-$1$ amplitude).
\end{enumerate}
The shadow obstruction tower controls the dependence of
$\int_{\Sigma_g} \cA$ on the moduli of $\Sigma_g$.
\end{proposition}

\begin{proposition}[Genus tower from FH]
\label{prop:fh-genus-tower}
\ClaimStatusConditional
For a CY$_2$ category $\cC$ with $\Etwo$-chiral algebra
$\cA = \Phi(\cC)$, the factorization homology genus tower satisfies
\[
 \chi\!\left(\int_{\Sigma_g} \cA\right)
 \;=\;
 \kappa_{\mathrm{cat}}(\cC) \cdot \lambda_g
 \qquad\textup{(UNIFORM-WEIGHT)}
\]
at genus $g \geq 1$, where $\kappa_{\mathrm{cat}} = \chi^{\CY}(\cC)$
is the CY Euler characteristic
(Theorem~\textup{\ref{thm:cy-modular-characteristic}}). For
multi-weight algebras at $g \geq 2$, the scalar formula receives
the cross-channel correction $\delta F_g^{\mathrm{cross}}$ of Vol~I,
Theorem~D. The conditional dependence is on Theorem~CY-B(iii)
(curvature identification).
\end{proposition}

\begin{remark}[FH versus chiral homology]
\label{rem:fh-vs-chiral}
The Ayala--Francis factorization homology $\int_{\Sigma_g} \cA$ and the
Beilinson--Drinfeld chiral homology $H^{\mathrm{ch}}_\bullet(\Sigma_g, \cA)$
are equivalent for $\Etwo$-algebras that arise as factorization envelopes
of chiral Lie algebras (by the comparison theorem of Francis, 2013).
In the CY setting, $\cA = \Phi(\cC)$ is always of this type when CY-A
holds (at $d = 2$), so the two invariants agree. The FH formulation has
the advantage of applying to topological $\Etwo$-algebras without
holomorphic structure.
\end{remark}


\section{The $E_2$-equivariant bar complex}
\label{sec:e2-equivariant-bar}

%% FOUR OBJECTS:
%% 1. B(A) = bar coalgebra = T^c(s^{-1} A-bar) with deconcatenation + twist
%% 2. A^i = H^*(B(A)) = dual coalgebra (Koszul cohomology of bar)
%% 3. A^! = ((A^i)^v) = dual ALGEBRA (linear or Verdier dual)
%% 4. Z^der_ch(A) = derived chiral center = Hochschild cochains = bulk

The $E_1$-bar complex $B^{\mathrm{ord}}(\cA) = T^c(s^{-1}\bar{\cA})$
of Volume~I is an ordered coalgebra with deconcatenation coproduct and
a single differential from OPE residues along one real direction.
In the $\Etwo$ setting, the two $\Eone$-directions of the Dunn
decomposition $\Etwo \simeq \Eone \otimes_{E_0} \Eone$ each contribute
a bar differential. The $\Etwo$-equivariant bar complex intertwines them
through the $R$-matrix.

\begin{definition}[$E_2$-equivariant bar complex]
\label{def:e2-equivariant-bar}
For an $\Etwo$-chiral algebra $\cA$ with augmentation ideal
$\bar{\cA} = \ker(\varepsilon)$, the \emph{$E_2$-equivariant bar
complex} is the graded coalgebra
\[
 B_{E_2}(\cA)
 \;=\;
 T^c(s^{-1}\bar{\cA})
\]
equipped with:
\begin{enumerate}[label=(\roman*)]
 \item two commuting differentials $d_X, d_Y$ from OPE residues
 in each $\Eone$-direction, satisfying $d_X^2 = d_Y^2 = 0$ and
 $d_X \circ d_Y = d_Y \circ d_X$;
 \item two deconcatenation coproducts $\Delta_X, \Delta_Y$ giving
 $B_{E_2}(\cA)$ the structure of an $\Etwo$-coalgebra;
 \item a braid group action: the symmetric group $S_n$ action on
 $B^{\mathrm{ord}}_n(\cA) = (s^{-1}\bar{\cA})^{\otimes n}$ lifts
 to a braid group $B_n$ action via the $R$-matrix
 $\cR_\cA(z)$ (Definition~\ref{def:categorical-r-matrix}).
\end{enumerate}
The total differential is $d = d_X + d_Y$; it satisfies $d^2 = 0$ on
the nose (genus-$0$ bar complex). At genus $g \geq 1$, the fiberwise
differential satisfies $d_{\mathrm{fib}}^2 = \kappa_{\mathrm{cat}} \cdot \omega_g$
(Hodge curvature; see Theorem~\ref{thm:e2-bar-cobar}(iii)).
\end{definition}

\begin{remark}[$B_{E_2}(\cA)$ is not a Swiss-cheese coalgebra]
\label{rem:e2-bar-not-sc}
The $\Etwo$-bar complex $B_{E_2}(\cA)$ is an $\Etwo$-coalgebra, \emph{not}
a Swiss-cheese ($\SC^{\mathrm{ch,top}}$) coalgebra. The SC structure of
Vol~II is two-coloured (holomorphic + topological) and emerges on the
\emph{derived center} $\Zder(\cA) = C^\bullet_{\mathrm{ch}}(\cA, \cA)$
(the bulk algebra), not on the bar complex of $\cA$ itself. The bar
complex encodes the boundary; the SC structure governs the bulk-boundary
coupling. Concretely:
\begin{itemize}
 \item $B_{E_2}(\cA)$: $\Etwo$-coalgebra, boundary data,
 deconcatenation coproduct, braid group action;
 \item $\Zder(\cA)$: $\Etwo$-algebra (by the BZF theorem,
 Theorem~\ref{thm:bzfn}), bulk data, carries the SC structure
 when coupled with the boundary $\cA$.
\end{itemize}
The identification $\cZ(\Rep^{E_1}(\cA)) \simeq \Rep^{E_2}(\Zder(\cA))$
(Chapter~\ref{ch:drinfeld-center}) relates the two: the braided monoidal
structure on representations of the bulk is the Drinfeld center of the
monoidal structure on representations of the boundary.
\end{remark}

\begin{proposition}[Three bar complexes and their relation]
\label{prop:three-bars-relation}
\ClaimStatusProvedElsewhere
For an $\Etwo$-chiral algebra $\cA$, the three bar complexes of
Construction~\ref{constr:three-bar-complexes} are related by:
\[
 B_{E_\infty}^n(\cA) = B_{E_1}^n(\cA) / S_n, \qquad
 H^*(B_{E_\infty}^n(\cA)) = H^*(B_{E_2}^n(\cA))^{S_n}.
\]
The $\Eone$-bar is the primitive ordered object; the $\Etwo$-bar
refines the $S_n$-action to the braid group $B_n$ via the $R$-matrix;
the $\Einf$-bar is the $S_n$-coinvariant shadow. The averaging map
$\mathrm{av} \colon B_{E_1} \to B_{E_\infty}$ factors through $B_{E_2}$:
\[
 B_{E_1} \xrightarrow{B_n\text{-equivariantization}} B_{E_2}
 \xrightarrow{B_n \twoheadrightarrow S_n} B_{E_\infty}.
\]
\end{proposition}


\section{The braided bar-cobar adjunction}
\label{sec:braided-bar-cobar}

The Vol~I bar-cobar adjunction (Theorems~A and~B) operates at the $\Eone$
level: $B^{\mathrm{ord}}$ is an $\Eone$-coalgebra and $\Omega$ recovers
the $\Eone$-algebra. The $\Etwo$-refinement lifts the adjunction to
braided coalgebras, with the $R$-matrix providing the additional datum
that distinguishes $\Etwo$ from $\Eone$.

\begin{theorem}[$E_2$-bar-cobar adjunction: Theorem CY-B]
\label{thm:e2-bar-cobar}
\ClaimStatusConditional
For an $\Etwo$-chiral algebra $\cA$, the bar-cobar adjunction
\[
 B_{E_2} \colon E_2\text{-}\mathrm{ChirAlg}
 \rightleftarrows E_2\text{-}\mathrm{ChirCoalg}
 \cocolon \Omega_{E_2}
\]
satisfies:
\begin{enumerate}[label=(\roman*)]
 \item The $\Etwo$ bar complex $B_{E_2}(\cA)$ is a factorization
 $\Etwo$-coalgebra on $\Ran(X) \times \Ran(Y)$;
 \item Cobar-bar inversion:
 $\Omega_{E_2}(B_{E_2}(\cA)) \simeq \cA$ on the Koszul locus;
 \item The CY trace manifests as curvature: at genus $g \geq 1$,
 $d_{\mathrm{fib}}^2 = \kappa_{\mathrm{cat}} \cdot \omega_g$,
 where $\omega_g = c_1(\lambda)$ is the Hodge class on
 $\overline{\cM}_g$ and $\kappa_{\mathrm{cat}} = \chi^{\CY}(\cC)$.
\end{enumerate}
\end{theorem}

\begin{remark}[Status of Theorem CY-B]
\label{rem:cy-b-status}
Items~(i) and~(ii) follow from the Volume~I bar-cobar machine
(Theorems~A and~B) applied to each $\Eone$-direction separately, with
compatibility ensured by Dunn additivity
(Proposition~\ref{prop:dunn-e2}). The argument proceeds in three steps.

\emph{Step 1: $\Eone \times \Eone$ bar-cobar.} Apply the Vol~I
bar-cobar adjunction to each $\Eone$-factor independently. This produces
two commuting bar-cobar adjunctions $(B_X, \Omega_X)$ and
$(B_Y, \Omega_Y)$.

\emph{Step 2: $R$-matrix intertwining.} The $R$-matrix
$\cR_\cA(z) \in \End(V \otimes V)[[z, z^{-1}]]$ at arity $(1,1)$
intertwines the two bar differentials: $\cR \circ d_X = d_Y \circ \cR$
on the arity-$(1,1)$ component. This is the braided compatibility
condition.

\emph{Step 3: Higher coherences.} The full $\Etwo$-equivariant
refinement requires coherent intertwining at all arities, governed by
the braid group $B_n$ action. Steps~1 and~2 are proved; the detailed
verification of higher coherences (arity $\geq 4$) is established as a
rigorous proof sketch. The result is unconditional at arity $\leq 3$
and conditional on the higher-arity coherences beyond that.

Item~(iii) uses Theorem~D of Volume~I with the CY Euler
characteristic $\kappa_{\mathrm{cat}} = \chi^{\CY}(\cC)$
(Theorem~\ref{thm:cy-modular-characteristic}). The curvature
$d_{\mathrm{fib}}^2 = \kappa_{\mathrm{cat}} \cdot \omega_g$ is the
\emph{fiberwise} statement on $\overline{\cM}_g$; the bar differential
$d^2_{\mathrm{bar}} = 0$ always holds (the curvature is Hodge, not bar).
\end{remark}

\begin{proposition}[Arity-$(1,1)$ of $B_{E_2}$]
\label{prop:e2-bar-arity-11}
\ClaimStatusProvedHere
The arity-$(1,1)$ component of $B_{E_2}(\cA)$ is the space
$\End(V \otimes V)$ equipped with the $R$-matrix
$\cR_\cA(z)$. The bar differential restricted to this component
recovers the unitarity condition $\cR_{12}(z) \cR_{21}(-z) = \id$
(strong unitarity), and the braiding constraint
$\cR_{12}(z) = \tau \circ \sigma_{12}(z)$.
\end{proposition}

\begin{proof}
The $E_2$-bar complex at arity $(1,1)$ consists of elements
$\alpha \in (s^{-1}\bar{\cA})^{\otimes 2}$ with one element
from each $E_1$-direction. The bar differential
$d(\alpha) = \mu(\alpha)$ is the OPE residue extracted via
$d\log(z_1 - z_2)$ (Volume~I, bar differential). The
compatibility of the two $E_1$-structures forces
$d_{X} \circ d_{Y} = d_{Y} \circ d_{X}$; evaluating on
$(s^{-1}v) \otimes (s^{-1}w)$ gives
$\cR_{12}(z)\cR_{21}(-z) = \id$ (strong unitarity).
The braiding $\sigma_{12}(z) = \cR_{12}(z) \circ \tau^{-1}$
follows from the identification of $\cR$ with the
half-braiding (Definition~\ref{def:categorical-r-matrix}).
\end{proof}


\section{Quantum group structure from braided Koszul duality}
\label{sec:qg-from-braided-koszul}

The Vol~I Koszul dual $A^! = (H^*(B(A))^{\vee})$ is a linear-dual
algebra. In the $\Etwo$ setting, the bar cohomology
$H^*(B_{E_2}(\cA))$ carries a quasi-triangular Hopf algebra structure
from the $R$-matrix, and the question becomes: does this Hopf algebra
coincide with the quantum group $\Uq(\frakg)$ when the input category
has $\frakg$-symmetry?

\begin{conjecture}[Braided bar cohomology recovers $\Uq(\frakg)$]
\label{conj:braided-bar-uq}
\ClaimStatusConjectured
For the $\Etwo$-chiral algebra $\cA = \Phi(\cC)$ associated to a
CY$_2$ category $\cC$ with $\frakg$-symmetry, the braided bar
cohomology
\[
 H^*(B_{E_2}(\cA)) \;\simeq\; \Uq(\frakg)
\]
as quasi-triangular Hopf algebras, where $q = e^{2\pi i / (k + h^\vee)}$
and $k$ is the level determined by the CY Euler characteristic
$\kappa_{\mathrm{cat}}(\cC)$.
\end{conjecture}

\begin{remark}[Evidence and scope]
\label{rem:braided-bar-evidence}
Conjecture~\ref{conj:braided-bar-uq} is verified in two cases:
\begin{enumerate}[label=(\roman*)]
 \item For $\cC = D^b(\mathrm{Coh}(\mathrm{pt}/G))$ with $G$ simple
 and simply connected, the CY-to-chiral functor produces
 $V_k(\frakg)$ (by CY-A at $d = 2$). The $\Etwo$-bar cohomology
 should recover $\Uq(\frakg)$ via the Kazhdan--Lusztig equivalence
 (Theorem~\ref{thm:qgf-kazhdan-lusztig}). This is verified at the
 level of categories: $\Rep_q(\frakg) \simeq \cO_{k}(\hat{\frakg})$
 at roots of unity.
 \item For $\cC = D^b(\mathrm{Coh}(K3))$, the chiral algebra is a lattice
 VOA and the $\Etwo$-bar cohomology recovers a lattice quantum group
 (Example~\ref{ex:braided-k3}).
\end{enumerate}
The conjecture is open for non-geometric CY categories and for $d = 3$
(where CY-A$_3$ is itself conjectural).
\end{remark}

\begin{remark}[Three dualities]
\label{rem:braided-koszul-langlands}
%: Three distinct operations must be kept separate.
The braided Koszul dual $\cA^!_{E_2}$ produces the quantum group
$\Uq(\frakg)$. This must be distinguished from:
\begin{enumerate}[label=(\alph*)]
 \item the Feigin--Frenkel involution $k \mapsto -k - 2h^\vee$, which
 acts within the same family of affine algebras;
 \item the Langlands dual $\frakg \mapsto \frakg^L$ (root-coroot
 exchange), which produces a genuinely different algebra.
\end{enumerate}
For $\fsl_2$ (self-Langlands-dual), all three operations coincide on
representation categories. For non-simply-laced types (e.g.,
$G_2, F_4$), they diverge. The braided Koszul dual exchanges the
braiding parameter $q \leftrightarrow q^{-1}$; the Langlands dual
exchanges the Cartan matrix $A \leftrightarrow A^T$; the Feigin--Frenkel
involution exchanges levels $k \leftrightarrow -k - 2h^\vee$.
\end{remark}


\section{The braided shadow obstruction tower}
\label{sec:braided-shadow-tower}

The Vol~I shadow obstruction tower $\Theta_A$ is the $\Sigma_n$-coinvariant
projection of a richer $E_1$ object. In the $\Etwo$ setting, the braided
shadow tower $\Theta^{E_2}_\cA$ sits between the ordered $E_1$ tower and
the symmetric $E_\infty$ shadow, with the braid group $B_n$ replacing
the symmetric group $S_n$.

\begin{definition}[Braided modular characteristic]
\label{def:braided-kappa}
For an $\Etwo$-chiral algebra $\cA$, the \emph{braided modular
characteristic} is the pair
\[
 (\kappa_{\mathrm{cat}}(\cA), \, \cR_\cA(z))
\]
consisting of the scalar modular characteristic
$\kappa_{\mathrm{cat}} = \mathrm{av}(\cR_\cA(z))$ (the
$\Sigma_2$-coinvariant of the $R$-matrix) and the full
$R$-matrix itself. The scalar $\kappa_{\mathrm{cat}}$ controls
the genus tower
$F_g = \kappa_{\mathrm{cat}} \cdot \lambda_g^{FP}$ (UNIFORM-WEIGHT);
the $R$-matrix controls the ordered/quantum-group sector.
\end{definition}

\begin{proposition}[Braided shadow structure]
\label{prop:braided-shadow}
\ClaimStatusProvedHere
The $\Etwo$-bar complex $B_{E_2}(\cA)$ carries a braided shadow
obstruction tower $\Theta^{E_2}_\cA$ whose arity-by-arity
projection under the averaging map recovers the Volume~I
shadow obstruction tower:
\[
 \mathrm{av}(\Theta^{E_2}_{\cA, n})
 = \Theta_{A_\cC, n}
 \quad \text{for all $n \geq 2$.}
\]
At arity $2$: $\mathrm{av}(\cR(z)) = \kappa_{\mathrm{cat}}$.
At arity $3$: $\mathrm{av}(\Phi_{KZ}) = C$ (the cubic shadow of
Volume~I).
\end{proposition}

\begin{remark}[The three shadow tiers]
\label{rem:three-shadow-tiers}
The three bar complexes produce three tiers of the shadow tower:
\begin{center}
\renewcommand{\arraystretch}{1.2}
\begin{tabular}{lll}
 \toprule
 Bar complex & Group action & Shadow data \\
 \midrule
 $B_{E_1}$ (ordered) & trivial ($1$) & full $R$-matrix $\cR(z)$, Yangian \\
 $B_{E_2}$ (braided) & braid group $B_n$ & $B_n$-equivariant tower, quantum group \\
 $B_{E_\infty}$ (symmetric) & symmetric $S_n$ & scalar $\kappa_{\mathrm{cat}}$, genus tower \\
 \bottomrule
\end{tabular}
\end{center}
Each tier retains strictly less information than the one above.
The averaging map $\mathrm{av}$ is the $B_n \twoheadrightarrow S_n$
projection applied to each arity. It is lossy: the $R$-matrix is
forgotten, and only its $S_2$-coinvariant $\kappa_{\mathrm{cat}}$
survives.
\end{remark}

\begin{example}[Braided structure for $K3$]
\label{ex:braided-k3}
For a K3 surface $S$ with $\cC = D^b(\Coh(S))$, the
$\Etwo$-chiral algebra has $\kappa_{\mathrm{cat}} = 2$. The
$R$-matrix is the Casimir
$\cR(z) = 1 + \kappa_{\mathrm{cat}}\,\Omega/z + O(z^{-2})$,
where $\Omega$ is the Mukai pairing on $H^*(S, \Z)$. The braided
bar cohomology $H^*(B_{E_2}(\Phi(\cC)))$ recovers the lattice
quantum group associated to the Mukai lattice
$\Lambda_{K3} = U^3 \oplus E_8(-1)^2$.
The shadow class is $\mathbf{G}$ (Gaussian, $r_{\max} = 2$):
the lattice VOA is a free-field algebra and the shadow tower
terminates at arity $2$ (both at the $E_1$ and $E_2$ levels).
\end{example}


\section{Braided complementarity}
\label{sec:braided-complementarity}

Vol~I Theorem~C establishes complementarity:
$\kappa_{\mathrm{ch}}(A) + \kappa_{\mathrm{ch}}(A^!) = K$
(the Koszul conductor, family-dependent). The $\Etwo$-refinement
upgrades this scalar relation to an $R$-matrix-level statement.

\begin{conjecture}[Braided complementarity]
\label{conj:braided-complementarity}
\ClaimStatusConjectured
For an $\Etwo$-chiral algebra $\cA$ on the Koszul locus, the braided
Koszul dual $\cA^!_{E_2}$ satisfies:
\begin{enumerate}[label=(\roman*)]
 \item \emph{Scalar complementarity}:
 $\kappa_{\mathrm{cat}}(\cA) + \kappa_{\mathrm{cat}}(\cA^!_{E_2}) = K$,
 the family-dependent Koszul conductor (Vol~I, Theorem~C);
 \item \emph{$R$-matrix inversion}:
 $\cR_{\cA^!_{E_2}}(z) = \cR_\cA(z)^{-1}$ as formal power series
 in $z^{-1}$, so the braided Koszul dual reverses the braiding;
 \item \emph{Conductor from $R$-matrix}:
 the Koszul conductor $K$ is recovered from the $R$-matrix
 as $K = \mathrm{av}(\cR_\cA(z)) + \mathrm{av}(\cR_\cA(z)^{-1})$.
\end{enumerate}
\end{conjecture}

\begin{remark}[Status of braided complementarity]
\label{rem:braided-complementarity-status}
Item~(i) is an immediate consequence of Vol~I Theorem~C applied
to the underlying $\Einf$-data. Item~(ii) is the genuine $\Etwo$
content: it asserts that Koszul duality acts on the $R$-matrix by
inversion, so $q \mapsto q^{-1}$. For the Kazhdan--Lusztig
equivalence, this corresponds to the duality
$\Rep_q(\frakg) \simeq \Rep_{q^{-1}}(\frakg)^{\mathrm{rev}}$,
which is the standard anti-braiding equivalence for quantum groups.
Item~(iii) follows formally from (i) and (ii), since
$\mathrm{av}(\cR^{-1}) = \kappa_{\mathrm{cat}}(\cA^!_{E_2})$.
The conjecture is conditional on Theorem~CY-B (the braided bar-cobar
adjunction at the $\Etwo$ level) and on
Conjecture~\ref{conj:braided-bar-uq} (braided bar cohomology
recovering the quantum group).
\end{remark}
